% from ega0-21-fr.tex
\oldpage[0\textsubscript{IV-1}]{180}
  Les formules (21.6.4.4) et (21.6.4.2) donnent d'autre part
  \[
    d(E/M_0,M)=d(E/L_0,M)+d(L_0/M_0,M),
  \]
  \[
    d(L_0/M_0,M)=d(L_0/M_0,M/N)+d(L_0/M_0,N),
  \]
  d'où
  \[
    \delta=d(E/L_0,M)+d(L_0/M_0,N)-d(E/L_0,k_0)
    +d(L_0/M_0,M/N)-r.
  \]

  Nous allons montrer que l'on a la relation
  \begin{equation}
  \tag{21.9.6.5}
  \label{0.21.9.6.5-fr}
    d(L_0/M_0,M/N)=r.
  \end{equation}
  Pour cela, remarquons que, par définition, on a
  \[
    d(L_0/M_0,M/N)
    =\operatorname{rg}_{L_0}\Upsilon_{L_0/M/N}
    -\operatorname{rg}_{M_0}\Upsilon_{M_0/M/N}.
  \]
  Compte tenu des deux suites exactes
  \[
  \begin{aligned}
    0&\longrightarrow\Upsilon_{L_0/M/N}
      \longrightarrow\Omega^1_{M/N}\otimes_ML_0
      \longrightarrow\Omega^1_{L_0/N}
      \longrightarrow\Omega^1_{L_0/M}\longrightarrow0,\\
    0&\longrightarrow\Upsilon_{M_0/M/N}
      \longrightarrow\Omega^1_{M/N}\otimes_MM_0
      \longrightarrow\Omega^1_{M_0/N}
      \longrightarrow\Omega^1_{M_0/M}\longrightarrow0,
  \end{aligned}
  \]
  il vient
  \[
  \begin{split}
    d(L_0/M_0,M/N)={}&\operatorname{rg}_{L_0}\Omega^1_{L_0/M}
      -\operatorname{rg}_{L_0}\Omega^1_{L_0/N}\\
      &-\operatorname{rg}_{M_0}\Omega^1_{M_0/M}
      +\operatorname{rg}_{M_0}\Omega^1_{M_0/N}.
  \end{split}
  \]
  Or, on a $N(L_0^p)\supset M$, donc
  $\Omega^1_{L_0/N}=\Omega^1_{L_0/M}$ (21.1.5, (ii)). D'autre part, les
  $T_i^p$ $(1\leqslant i\leqslant r)$ forment une $p$-base de $M$ sur $N$,
  et une $p$-base de $K_0$ sur $k$ est aussi une $p$-base de $M_0$ sur $M$,
  donc, comme $M_0^p\subset N$,
  \[
    \operatorname{rg}_{M_0}\Omega^1_{M_0/N}
    =r+\operatorname{rg}_{K_0}\Omega^1_{K_0/k}
    =r+\operatorname{rg}_{M_0}\Omega^1_{M_0/M}
  \]
  en vertu de (21.4.5) et (21.1.10), ce qui prouve (21.9.6.5).

  On a encore, par (21.6.4.2),
  \[
    d(E/L_0,M)=d(E/L_0,k_0)+d(E/L_0,M/k_0)
  \]
  et
  \[
    d(L_0/M_0,N)=d(L_0/M_0,N/k)+d(L_0/M_0,k).
  \]
  Mais comme $L_0$ est séparable sur $K_0$ (21.9.6.4), on a
  $d(L_0/K_0,k)=0$, et par suite aussi $d(L_0/M_0,k)=0$ ((21.6.2) et
  (21.6.4)). Il vient donc
  \[
    \delta=d(E/L_0,M/k_0)+d(L_0/M_0,N/k).
  \]
  Appliquons encore deux fois (21.6.4.2) au premier terme du second membre de
  cette formule~; on obtient
  \[
    d(E/L_0,M/k_0)=d(E/L_0,M/N)+d(E/L_0,N/k)+d(E/L_0,k/k_0).
  \]
  Comme $M\subset N(L_0^p)$, on a $d(E/L_0,M/N)=0$ (21.6.2), donc (par
  (21.6.4.4))
  \[
    \delta=d(E/L_0,k/k_0)+d(E/M_0,N/k).
  \]
  Mais $M_0$ et $L_0$ sont des extensions séparables de $K_0$ (21.9.6.4),
  donc
  \[
    d(L_0/K_0,k/k_0)=d(M_0/K_0,k/k_0)=0
  \]
  (21.6.2), et appliquant encore deux fois (21.6.4.4),

  